\subsection{Conclusion}
This thesis has shown that the proposed Side Scan Sonar SLAM data processing pipeline can run in real time on real embedded hardware using real sensor data from realistic marine conditions. The system is able to process raw side scan sonar data, correct and normalize sonar swaths, generate stable local maps, extract repeatable landmarks, estimate approximate landmark height, assign useful descriptors, and publish the required outputs while staying within the runtime limits of a Raspberry Pi 4 class platform. The results show that the pipeline is not only fast, but also produces geometrically consistent maps and reliable feature representations under real seabed conditions. This demonstrates that the pipeline is suitable both for offline analysis and as a practical foundation for onboard Side Scan Sonar SLAM on modern autonomous underwater vehicles.
\\ \\
A key result is that the Side Scan Sonar SLAM data processing pipeline is highly tunable and scalable. Parameters such as map resolution, pruning thresholds, chunk configuration, update rate, and feature extraction settings allow the pipeline to be adapted to different hardware limits and mission requirements. The results show that map resolution and memory usage are especially important for runtime, and that the system can be tuned to reduce computational load while still preserving useful map and landmark information. This makes the approach practical for constrained systems where perception must share compute resources with state estimation, control, guidance, logging, and communication.
\\ \\
The work has also produced a clean and reusable software foundation. Significant effort has been put into making the code modular, documented, configurable, and easy to integrate. For ROS2 based vehicles, the system can be deployed directly as ROS2 packages with exposed parameters and ready to use interfaces. For custom vehicles or systems that do not use ROS2, the core processing is available as general reusable libraries without depending on the ROS2 wrappers. This makes the implementation portable across different drone platforms and reduces the amount of work needed for future integration.
\\ \\
Although the full closed loop SLAM system was not completed, this work establishes the real-time embedded foundation needed to make it achievable. The Side Scan Sonar SLAM data processing pipeline now provides stable local maps, compact landmark representations, reusable software libraries, and practical ROS2 integration that future work can build directly upon. By moving Side Scan Sonar SLAM data processing from offline experiments toward real-time operation on real autonomous vehicle hardware, this thesis helps pave the way for more robust underwater navigation, longer autonomous missions, and future SLAM research using side scan sonar as a practical onboard perception sensor.
